\section{An exact structural bridge: annealable rank-factor layers}
The general spectral theory still has a structural-realizability problem: a prediction eigenvector is not a literal neuron or matrix entry. A useful architecture should make an irreversible rank contraction an exact coordinate operation rather than an approximate analogy. One simple construction is an additive rank-factor layer.

Consider an affine layer
\[
z=W h+b,
\qquad
W=\sum_{q=1}^{r}s_q u_qv_q^\top,
\]
where $u_q\in\mathbb R^{d_{\rm out}}$, $v_q\in\mathbb R^{d_{\rm in}}$, and $s_q\in\mathbb R$ are trainable. No orthogonality is assumed here; the $q$'s are structural rank-one components, not automatically singular vectors. Let the remaining network parameters be denoted by $\xi$.

The raw factorization has a gauge symmetry that must be fixed before scale-based structural diagnostics are interpreted. For every nonzero $c,d$,
\[
(s_q,u_q,v_q)
\mapsto
\left(\frac{s_q}{cd},cu_q,dv_q\right)
\]
leaves $s_qu_qv_q^\top$ unchanged. Hence raw $|s_q|$ and Euclidean parameter-gradient norms are not representation invariant.

\begin{proposition}[Canonical component gauge]\label{prop:rank-gauge}
Every nondegenerate rank-one component admits the function-preserving canonicalization
\[
\bar u_q=\frac{u_q}{\norm{u_q}},
\qquad
\bar v_q=\frac{v_q}{\norm{v_q}},
\qquad
\bar s_q=s_q\norm{u_q}\norm{v_q}.
\]
Then
\[
\norm{\bar u_q}=\norm{\bar v_q}=1,
\qquad
|\bar s_q|=\norm{s_qu_qv_q^\top}_{F}.
\]
Thus the magnitude of the canonical scale is the gauge-invariant Frobenius magnitude of the represented component. Canonicalization changes neither the layer matrix nor the network predictor.
\end{proposition}

\begin{proof}
Substitution gives $\bar s_q\bar u_q\bar v_q^\top=s_qu_qv_q^\top$. The Frobenius norm of an outer product is the product of the vector norms, which gives the second identity.
\end{proof}

All scale and group-gradient diagnostics below are understood in this canonical gauge. A practical implementation can canonicalize a copied checkpoint before evaluating a candidate and reset the optimizer in both matched dense and reduced counterfactual continuations. This avoids attributing scientific meaning to an arbitrary factor scaling.

For a kept set $K\subset\{1,\ldots,r\}$ and dropped set $D=K^c$, define the physical compact layer
\[
W_K=\sum_{q\in K}s_qu_qv_q^\top.
\]
In the full parameter space define the \emph{ghost embedding} by leaving $u_q,v_q$ unchanged for $q\in D$ but setting $s_q=0$. Let $S_K$ be the coordinate projector that keeps $\xi,b$ and all parameters $(s_q,u_q,v_q)$ for $q\in K$ while freezing every parameter of the dropped groups.

\begin{theorem}[Exact physical embedding of rank contraction]\label{thm:rank-factor-embedding}
For every checkpoint and every kept set $K$:
\begin{enumerate}[label=(\roman*)]
\item the predictor of the physically compact rank-$|K|$ layer is exactly equal to the predictor of its ghost embedding in the rank-$r$ parameter space;
\item the projected gradient flow of the ghost model under $S_K$ is exactly the embedded gradient flow of the physically compact model;
\item if the compact initialization copies every kept/shared parameter and only zeros the dropped scales, then the full-space contraction jump is
\[
\boxed{
d_0=\norm{s_D}_2;}
\]
in the chosen gauge; in the canonical gauge each $|s_q|$ is also the component Frobenius magnitude;
\item deleting one component removes exactly $d_{\rm in}+d_{\rm out}+1$ trainable factor parameters and removes its two rank-factor matrix-vector contractions from the physical forward/backward graph.
\end{enumerate}
\end{theorem}

\begin{proof}
When $s_q=0$, the rank-one matrix $s_qu_qv_q^\top$ is identically zero, so removing the corresponding term or retaining it as a zero-scale ghost gives the same affine map and therefore the same network predictor. Derivatives with respect to every shared and kept parameter are also identical in the two representations. The projected ghost flow freezes the deleted coordinates, including the scale directions that could otherwise regrow, so its active-coordinate ODE is exactly the compact-model ODE. The only coordinates changed at the embedding checkpoint are $s_q$, $q\in D$, giving the Euclidean distance $\norm{s_D}_2$. The parameter-count statement follows directly from the shapes of $u_q,s_q,v_q$.
\end{proof}

This construction solves the structural bridge for one concrete family: rank can be reduced during training with a literal decrease in parameters and arithmetic. It does \emph{not} say that small $|s_q|$ alone makes a component useless. Target value is still required.

At a canonicalized checkpoint define the exact structural group-gradient energy
\[
E_D^{\rm grp}(t)
:=\norm{(I-S_K)g_t}^2
=
\sum_{q\in D}
\left[
\norm{\nabla_{u_q}\R_t}^2
+(\partial_{s_q}\R_t)^2
+\norm{\nabla_{v_q}\R_t}^2
\right].
\]
Because the groups occupy orthogonal parameter coordinates, this identity is exact in the fixed canonical coordinate representation. It is a structural analogue of the target-weighted modal energy $\mu_ja_j^2$: it measures how much present loss-dissipation budget actually flows through the parameters that would disappear.

\begin{theorem}[Gap-free checkpoint certificate for physical parameter groups]\label{thm:group-certificate}
Let $S$ be any fixed coordinate/group projector representing a physically realizable contraction, and suppose on a common tube containing the dense and projected trajectories that $g$ is $L_g$-Lipschitz and $F$ is $B_{J,\mathrm{tube}}$-Lipschitz. Define
\[
E_S(t):=\norm{(I-S)g_t}^2,
\qquad
A_t(h):=\int_t^{t+h}\norm{g_s}\,ds.
\]
As in Theorem~\ref{thm:checkpoint-certificate}, set
\[
\overline A_t(H)
=\min\left\{
\sqrt{H\R_t},
\norm{g_t}\phi(L_g,H)
\right\}.
\]
For the integral of the path budget define
\[
\overline B_t(H)
:=
\min\left\{
\frac23H^{3/2}\sqrt{\R_t},
\norm{g_t}\psi(L_g,H)
\right\},
\]
where
\[
\psi(L,H)
:=
\begin{cases}
\dfrac{e^{LH}-1-LH}{L^2},&L>0,\\[6pt]
H^2/2,&L=0.
\end{cases}
\]
Then the omitted dense gradient of the physical groups satisfies
\[
\boxed{
\int_t^{t+H}\norm{(I-S)g_s}\,ds
\le
H\sqrt{E_S(t)}+L_g\overline B_t(H).}
\]
If the contraction has initial full-space embedding error $d_0$, the terminal prediction discrepancy satisfies
\[
\boxed{
\norm{F(\theta_{t+H})-F(\widetilde\theta_{t+H})}
\le
B_{J,\mathrm{tube}}e^{L_gH}
\left[
d_0+H\sqrt{E_S(t)}+L_g\overline B_t(H)
\right].}
\]
For the canonically gauged rank-factor contraction of Theorem~\ref{thm:rank-factor-embedding}, one may take $d_0=\norm{s_D}_2$ and $E_S(t)=E_D^{\rm grp}(t)$ exactly.
\end{theorem}

\begin{proof}
Because $I-S$ is an orthogonal projector and $g$ is Lipschitz,
\[
\begin{aligned}
\norm{(I-S)g_s}
&\le
\norm{(I-S)g_t}
+\norm{(I-S)(g_s-g_t)}\\
&\le
\sqrt{E_S(t)}+L_g\norm{\theta_s-\theta_t}\\
&\le
\sqrt{E_S(t)}+L_gA_t(s-t).
\end{aligned}
\]
Integrating gives
\[
\int_t^{t+H}\norm{(I-S)g_s}ds
\le H\sqrt{E_S(t)}+L_g\int_0^H A_t(h)dh.
\]
The loss-dissipation path bound gives $A_t(h)\le\sqrt{h\R_t}$, whose integral is $\frac23H^{3/2}\sqrt{\R_t}$. The current-gradient Lipschitz bound gives
$A_t(h)\le\norm{g_t}\phi(L_g,h)$, whose integral is $\norm{g_t}\psi(L_g,H)$. Taking the smaller valid integral envelope gives $\overline B_t(H)$. The prediction statement is Theorem~\ref{thm:gronwall} followed by the common-tube predictor Lipschitz bound. The rank-factor specialization follows from Theorem~\ref{thm:rank-factor-embedding} and the orthogonal group-energy identity.
\end{proof}

\begin{remark}[Why this is useful]
The spectral checkpoint theorem and the structural-group theorem play different roles. The spectral theorem answers \emph{which capacity sector has little target-weighted future value and is unlikely to rotate back into relevance}. The group theorem answers \emph{whether the actual parameters we intend to delete are currently dynamically inactive enough to remove}. A practical rank-annealing controller should require both whenever both are available. This avoids the unjustified step ``small prediction eigenvalue $\Rightarrow$ delete arbitrary weights.''
\end{remark}

\begin{remark}[The scale is not a saliency score]
Even in canonical gauge, the term $\norm{s_D}_2$ controls the instantaneous embedding jump; it does not measure future usefulness. A component with tiny scale but large $|\partial_{s_q}\R|$ is exactly a component poised to regrow and should fail the group-gradient gate. Conversely a moderate scale can be removable after a function-preserving reparameterization or readout redistribution, but that requires a different embedding argument. The controller should therefore monitor canonical component magnitude, target-weighted group gradient, prediction-space discovery state, and finite-horizon risk jointly.
\end{remark}
